\subsection{Some Operations on Modules}

Let A and B be rings, and let P be an (A, B)-bimodule (II, §1, No.~14, p.~225). We will define two procedures, one to go from a left B-module to a left A-module, the other to go from a left A-module to a left B-module.

3.1. The operation \(\mathcal{T}\) . --- Let V be a left B-module. Denote the left A-module \(\mathrm{P} \otimes_{\mathrm{B}} \mathrm{V}\) (II, §3, No.~4, p.~247) by \(\mathcal{T}(\mathrm{V})\) . The law of action on \(\mathcal{T}(\mathrm{V})\) is given by the formula

\[
a (p \otimes v) = (a p) \otimes v\tag{1}
\]

for \(a \in \mathrm{A}\) , \(p \in \mathrm{P}\) , and \(v \in \mathrm{V}\) .

Let \(V'\) be a left B-module. For every B-linear mapping g from V to \(V'\) , the mapping \(1_{\mathrm{P}} \otimes g\) from \(\mathcal{T}(\mathrm{V})\) to \(\mathcal{T}(\mathrm{V}')\) is A-linear; we denote it by \(\mathcal{T}(g)\) . The mapping \(g \mapsto \mathcal{T}(g)\) from \(\operatorname{Hom}_{\mathrm{B}}(\mathrm{V}, \mathrm{V}')\) to \(\operatorname{Hom}_{\mathrm{A}}(\mathcal{T}(\mathrm{V}), \mathcal{T}(\mathrm{V}'))\) is Z-linear, and we have

\[
\mathcal {T} (1 _ {\mathrm{V}}) = 1 _ {\mathcal {T} (\mathrm{V})}, \quad \mathcal {T} (g ^ {\prime} \circ g) = \mathcal {T} (g ^ {\prime}) \circ \mathcal {T} (g)\tag{2}
\]

if V, \(V'\) , \(V''\) are left B-modules and \(g : V \to V'\) , \(g' : V' \to V''\) are B-linear mappings. Since the tensor product commutes with direct sums, if V is the direct sum of a family of submodules \((\mathrm{V}_{i})_{i \in \mathrm{I}}\) , then we can identify the A-module \(\mathcal{T}(\mathrm{V})\) with \(\bigoplus_{i} \mathcal{T}(\mathrm{V}_{i})\) .

3.2. The operation \(\mathcal{H}\) . --- Let M be a left A-module. Denote the left B-module \(\mathrm{Hom}_{\mathrm{A}}(\mathrm{P},\mathrm{M})\) (II, §1, No.~14, p.~225) by \(\mathcal{H}(\mathrm{M})\) . The law of action on \(\mathcal{H}(\mathrm{M})\) is given by the formula

\[
(b f) (p) = f (p b)\tag{3}
\]

for \(b \in \mathrm{B}\) , \(f \in \operatorname{Hom}_{\mathrm{A}}(\mathrm{P}, \mathrm{M})\) , and \(p \in \mathrm{P}\) .

Let \(M'\) be a left A-module. For every A-linear mapping g from M to \(M'\) , the mapping \(\mathrm{Hom}_{\mathrm{A}}(1_{\mathrm{P}}, g)\) from \(\mathcal{H}(\mathrm{M})\) to \(\mathcal{H}(\mathrm{M}')\) is B-linear; we denote it by \(\mathcal{H}(g)\) . The mapping \(g \mapsto \mathcal{H}(g)\) from \(\mathrm{Hom}_{\mathrm{A}}(\mathrm{M}, \mathrm{M}')\) to \(\mathrm{Hom}_{\mathrm{B}}(\mathcal{H}(\mathrm{M}), \mathcal{H}(\mathrm{M}'))\) is Z-linear, and we have

\[
\mathcal {H} (1 _ {\mathrm{M}}) = 1 _ {\mathcal {H} (\mathrm{M})}, \quad \mathcal {H} (g ^ {\prime} \circ g) = \mathcal {H} (g ^ {\prime}) \circ \mathcal {H} (g)\tag{4}
\]

if M, \(M'\) , \(M''\) are left A-modules and \(g : M \to M'\) , \(g' : M' \to M''\) are A-linear mappings. Suppose, moreover, that P is a finitely generated A-module; if M is the direct sum of a family of submodules \((\mathrm{M}_{i})_{i}\) , then we can identify \(\mathcal{H}(\mathrm{M})\) with \(\bigoplus_{i} \mathcal{H}(\mathrm{M}_{i})\) by VIII, p.~57.

3.3. Relations Between \(\mathcal{T}\) and \(\mathcal{H}\) . --- By Proposition 1 of II, §4, No.~1, p.~268, for every left A-module M and every left B-module V, there exists a unique group isomorphism

\[
\gamma : \mathrm{Hom} _ {\mathrm{A}} (\mathcal {T} (\mathrm{V}), \mathrm{M}) \longrightarrow \mathrm{Hom} _ {\mathrm{B}} (\mathrm{V}, \mathcal {H} (\mathrm{M}))\tag{5}
\]

characterized by the relation

\[
(\gamma (h) (v)) (p) = h (p \otimes v)\tag{6}
\]

for \(h \in \operatorname{Hom}_{\mathrm{A}}(\mathcal{T}(\mathrm{V}), \mathrm{M})\) , \(v \in V\) , and \(p \in P\) . The isomorphism \(\gamma\) is called the adjunction isomorphism.

Let M be a left A-module. The A-module \(\mathcal{T}(\mathcal{H}(M))\) is simply the A-module \(P \otimes_{B} \operatorname{Hom}_{A}(P, M)\) . By applying the above to the B-module \(\mathcal{H}(M)\) , we see that the mapping

\[
\alpha_ {\mathrm{M}} = \gamma^ {- 1} (\operatorname{Id} _ {\mathcal {H} (\mathrm{M})}): \mathcal {T} (\mathcal {H} (\mathrm{M})) \longrightarrow \mathrm{M}
\]

is the unique mapping satisfying

\[
\alpha_ {\mathrm{M}} (p \otimes f) = f (p)\tag{7}
\]

for \(p \in \mathrm{P}\) and \(f \in \mathrm{Hom}_{\mathrm{A}}(\mathrm{P}, \mathrm{M})\) . We say that \(\alpha_{\mathrm{M}}\) is the canonical A-linear mapping from \(\mathcal{T}(\mathcal{H}(\mathrm{M}))\) to M. For every A-linear mapping \(g: \mathrm{M} \to \mathrm{M}'\) , we have a commutative diagram

\[
\begin{array}{c} \mathcal {T} (\mathcal {H} (M)) \xrightarrow {\alpha_ {M}} M \\ \Biggl \downarrow \mathcal {T} (\mathcal {H} (g)) \Biggl \downarrow g \\ \mathcal {T} (\mathcal {H} (M ^ {\prime})) \xrightarrow {\alpha_ {M ^ {\prime}}} M ^ {\prime}. \end{array}\tag{I}
\]

The inverse

\[
\gamma^ {- 1}: \mathrm{Hom} _ {\mathrm{B}} (\mathrm{V}, \mathcal {H} (\mathrm{M})) \longrightarrow \mathrm{Hom} _ {\mathrm{A}} (\mathcal {T} (\mathrm{V}), \mathrm{M})
\]

of the adjunction isomorphism coincides with the mapping \(h \mapsto \alpha_{\mathrm{M}} \circ \mathcal{T}(h)\) . Indeed, by (6) and (7), we have the relations

\[
\gamma^ {- 1} (h) (p \otimes v) = (h (v)) (p) = \alpha_ {\mathrm{M}} (p \otimes h (v)) = \alpha_ {\mathrm{M}} \circ \mathcal {T} (h) (p \otimes v)
\]

for all \(h\in \mathrm{Hom}_{\mathrm{B}}(\mathrm{V},\mathcal{H}(\mathrm{M}))\) , \(v\in \mathrm{V}\) , and \(p\in \mathrm{P}\) .

Let V be a B-module. The B-module \(\mathcal{H}(\mathcal{T}(V))\) is simply the B-module \(\mathrm{Hom}_{\mathrm{A}}(\mathrm{P},\mathrm{P}\otimes_{\mathrm{B}}\mathrm{V})\) . By applying (5) to the A-module \(\mathcal{T}(V)\) , we see that the B-linear mapping \(\beta_{\mathrm{V}}=\gamma(\mathrm{Id}_{\mathcal{T}(V)})\) from V to \(\mathcal{H}(\mathcal{T}(V))\) is characterized by the relation

\[
\beta_ {\mathrm{V}} (v) (p) = p \otimes v\tag{8}
\]

for \(p \in P\) and \(v \in V\) . We call \(\beta_{V}\) the canonical B-linear mapping from V to \(\mathcal{H}(\mathcal{T}(V))\) . For every B-module \(V'\) and every B-linear mapping \(g : V \to V'\) , we have a commutative diagram

\[
\begin{array}{l} \mathrm{V} \xrightarrow {\beta_ {\mathrm{V}}} \mathcal {H} (\mathcal {T} (\mathrm{V})) \\ \Big \downarrow_ {g} \qquad \qquad \qquad \Big \downarrow_ {\mathcal {H} (\mathcal {T} (g))} \\ \mathrm{V} ^ {\prime} \xrightarrow {\beta_ {\mathrm{V} ^ {\prime}}} \mathcal {H} (\mathcal {T} (\mathrm{V} ^ {\prime})). \end{array}\tag{II}
\]

Note that the adjunction morphism (5) coincides with the mapping that sends u to \(\mathcal{H}(u)\circ\beta_{\mathrm{V}}\) . Indeed, from relations (6) and (8), we deduce the equalities

\[
(\gamma (u) (v)) (p) = u (p \otimes v) = u \circ (\beta_ {\mathrm{V}} (v)) (p)
\]

for all \(u \in \mathrm{Hom}_{\mathrm{A}}(\mathcal{T}(\mathrm{V}), \mathrm{M})\) , \(v \in \mathrm{V}\) , and \(p \in \mathrm{P}\) .

Remarks. --- 1) Let V and V' be B-modules. The adjunction isomorphism

\[
\gamma : \mathrm{Hom} _ {\mathrm{A}} (\mathcal {T} (\mathrm{V}), \mathcal {T} (\mathrm{V} ^ {\prime})) \longrightarrow \mathrm{Hom} _ {\mathrm{B}} (\mathrm{V}, \mathcal {H} (\mathcal {T} (\mathrm{V} ^ {\prime})))
\]

satisfies the relation \(\gamma (\mathcal{T}(f)) = \beta_{\mathrm{V}^{\prime}}\circ f\) for every \(f\in \mathrm{Hom}_{\mathrm{B}}(\mathrm{V},\mathrm{V}^{\prime})\) because

\[
(\gamma (\mathcal {T} (f)) (v)) (p) = \mathcal {T} (f) (p \otimes v) = p \otimes f (v) = (\beta_ {\mathrm{V} ^ {\prime}} \circ f) (v) (p).
\]

Let M and M' be A-modules; the inverse of the adjunction isomorphism

\[
\gamma^ {- 1}: \mathrm{Hom} _ {\mathrm{B}} (\mathcal {H} (\mathrm{M}), \mathcal {H} (\mathrm{M} ^ {\prime})) \longrightarrow \mathrm{Hom} _ {\mathrm{A}} (\mathcal {T} (\mathcal {H} (\mathrm{M})), \mathrm{M} ^ {\prime})
\]

satisfies the relation \(\gamma^{-1}(\mathcal{H}(u)) = u \circ \alpha_{\mathrm{M}}\) for every \(u \in \operatorname{Hom}_{\mathrm{B}}(\mathrm{M}, \mathrm{M}')\) . Indeed, we have the relations

\[
\gamma^ {- 1} (\mathcal {H} (u)) (p \otimes v) = (\mathcal {H} (u) (v)) (p) = u (v (p)) = u \circ \alpha_ {\mathrm{M}} (p \otimes v)
\]

for all \(u \in \mathrm{Hom}_{\mathrm{B}}(\mathrm{M}, \mathrm{M}')\) , \(v \in \mathcal{H}(\mathrm{M})\) , and \(p \in \mathrm{P}\) .

\begin{enumerate}
\def\labelenumi{\arabic{enumi})}
\setcounter{enumi}{1}
\tightlist
\item
  Let \(\mathrm{M}\) be a left A-module. The B-linear mappings
\end{enumerate}

\[
\beta_ {\mathcal {H} (\mathrm{M})}: \mathcal {H} (\mathrm{M}) \to \mathcal {H} (\mathcal {T} (\mathcal {H} (\mathrm{M})))
\]

\[
\mathcal {H} (\alpha_ {\mathrm{M}}): \mathcal {H} (\mathcal {T} (\mathcal {H} (\mathrm{M}))) \to \mathcal {H} (\mathrm{M})
\]

satisfy the relation \(\mathcal{H}(\alpha_{\mathrm{M}})\circ\beta_{\mathcal{H}(\mathrm{M})}=1_{\mathcal{H}(\mathrm{M})}\) . They are not bijective in general.

Let V be a left B-module. The A-linear mappings

\[
\mathcal {T} (\beta_ {\mathrm{V}}): \mathcal {T} (\mathrm{V}) \to \mathcal {T} (\mathcal {H} (\mathcal {T} (\mathrm{V}))) \quad \text { and } \quad \alpha_ {\mathcal {T} (\mathrm{V})}: \mathcal {T} (\mathcal {H} (\mathcal {T} (\mathrm{V}))) \to \mathcal {T} (\mathrm{V})
\]

satisfy the relation \(\alpha_{\mathcal{T}(\mathrm{V})}\circ \mathcal{T}(\beta_{\mathrm{V}}) = 1_{\mathcal{T}(\mathrm{V})}\) . They are not bijective in general.

\begin{enumerate}
\def\labelenumi{\arabic{enumi})}
\setcounter{enumi}{2}
\tightlist
\item
  Suppose that P is finitely generated as an A-module. Let M be the direct sum of a family \((\mathrm{M}_{i})_{i\in\mathrm{I}}\) of A-modules. The A-modules \(\mathcal{T}(\mathcal{H}(\mathrm{M}))\) and \(\bigoplus_{i}\mathcal{T}(\mathcal{H}(\mathrm{M}_{i}))\) are canonically isomorphic. When we identify them, \(\alpha_{\mathrm{M}}\) is identified with \(\bigoplus_{i}\alpha_{\mathrm{M}_{i}}\) . Likewise, let V be the direct sum of a family \((\mathrm{V}_j)_{j\in \mathrm{J}}\) of B-modules. The B-module \(\mathcal{H}(\mathcal{T}(\mathrm{V}))\) is identified with \(\bigoplus_{j}\mathcal{H}(\mathcal{T}(\mathrm{V}_{j}))\) , and the linear mapping \(\beta_{\mathrm{V}}\) with \(\bigoplus_{j}\beta_{\mathrm{V}_{j}}\) .
\end{enumerate}

